% Generated draft: requires mathematical and source review.
\documentclass[11pt]{article}
\usepackage[margin=25mm]{geometry}
\usepackage{fontspec}
\setmainfont{Latin Modern Roman}
\usepackage{amsmath,amssymb,mathtools,unicode-math}
\setmathfont{Latin Modern Math}
\usepackage{enumitem,needspace,xurl,hyperref}
\setlength{\parindent}{0pt}
\setlength{\parskip}{5pt}
\setlength{\emergencystretch}{3em}
\newcommand{\N}{\mathbb{N}}
\newcommand{\Z}{\mathbb{Z}}
\newcommand{\Q}{\mathbb{Q}}
\newcommand{\R}{\mathbb{R}}
\newcommand{\C}{\mathbb{C}}
\newcommand{\F}{\mathbb{F}}
\newcommand{\A}{\mathbb{A}}
\newcommand{\PP}{\mathbb{P}}
\newcommand{\E}{\mathbb{E}}
\newcommand{\eps}{\varepsilon}
\newcommand{\dd}{\,\mathrm{d}}
\begin{document}
\section*{UnsolvedMath ID: 30003606}
\textbf{Source problem code:} catalogue entry
\subsection{Definitions and mathematical statement}
Let \(k\) be a field. A standard graded \(k\)-algebra is a commutative finitely generated algebra \(R=\bigoplus_{n\geq 0}R_n\) with \(R_0=k\), generated over \(k\) by \(R_1\). Require \(R\) to be an integral domain. Put \(\mathfrak m=\bigoplus_{n>0}R_n\) and \(d=\dim R\). The Cohen--Macaulay condition means \(\operatorname{depth}R_{\mathfrak m}=d\); the Gorenstein condition means that \(R_{\mathfrak m}\) has finite injective dimension as a module over itself. The source specifies no restrictions on the field, so the questions below are formulated over a fixed but otherwise unspecified \(k\); the intended quantification over fields requires confirmation.

Define the Hilbert series and its \(h\)-polynomial by
\[
H_R(t)=\sum_{n\geq0}(\dim_k R_n)t^n
       =\frac{h_R(t)}{(1-t)^d},
\qquad
h_R(t)=\sum_{i=0}^{s}h_i t^i,\quad h_s\ne0.
\]
For the Cohen--Macaulay rings under consideration, the coefficients \(h_i\) are nonnegative integers. The polynomial is called unimodal if there exists \(j\in\{0,\ldots,s\}\) such that
\[
h_0\leq h_1\leq\cdots\leq h_j
\quad\text{and}\quad
h_j\geq h_{j+1}\geq\cdots\geq h_s.
\]
It is called log-concave if
\[
h_i^2\geq h_{i-1}h_{i+1}
\qquad(1\leq i\leq s-1).
\]
Thus endpoint inequalities are not required, and log-concavity is vacuous when \(s\leq1\).

\textbf{Principal question.} Is \(h_R(t)\) unimodal for every standard graded Gorenstein domain \(R\)?

\textbf{Related questions.} Is \(h_R(t)\) unimodal for every standard graded Cohen--Macaulay domain \(R\)? What happens when unimodality is replaced by log-concavity? Applying that replacement to both ring classes asks whether every standard graded Gorenstein domain, respectively every standard graded Cohen--Macaulay domain, has a log-concave \(h\)-polynomial. The wording of the source does not explicitly distinguish these two log-concavity variants, so this expansion requires source review.

\textbf{Supplied status qualification.} The record labels this compound problem partially solved. Its dated assessment attributes a disproof of the general Cohen--Macaulay log-concavity assertion to a Niesi--Robbiano example, while describing Gorenstein-domain unimodality as open in general. These are supplied literature assessments, not independently verified mathematical or current-status conclusions. They do not explicitly settle the other variants.
\subsection{Short English statement}
The principal question asks whether the numerator coefficients of the Hilbert series of every standard graded Gorenstein domain increase and then decrease. The record also asks about Cohen--Macaulay domains and about log-concavity. Its supplied dated assessment reports a counterexample to general Cohen--Macaulay log-concavity and leaves Gorenstein-domain unimodality open. Field restrictions and the precise scope of the log-concavity variants require review.
\subsection{Sources}
\begin{itemize}
\item Supplied local UnsolvedMath registry record, permanent ID \(30003606\), source problem code \texttt{catalogue entry}, title \emph{Unimodality of Gorenstein Hilbert Series}; registry location \texttt{data/unsolvedmath/problems.json}, with \texttt{catalog/problems.json} designated as the repository's supplementary registry. This draft uses the supplied complete record; registry files were not independently inspected.
\item Supplied report citation: \emph{Mini-Workshop: Lattice Polytopes: Methods, Advances, Applications} (2018), Oberwolfach Reports, pp.\ 2659--2701, DOI \texttt{10.4171/owr/2017/44}. Supplied source URL: \url{https://doi.org/10.4171/owr/2017/44}. The record separately gives proposed year 2017 and names no individual proposer.
\item Supplied literature reference: Niesi--Robbiano, ``disproof of Hibi's conjecture,'' \url{https://arxiv.org/abs/alg-geom/9202020}. The record attributes the Cohen--Macaulay log-concavity counterexample to this reference.
\item Supplied literature reference: Conca--Singh--Varbaro, ``Gorenstein h-vectors,'' \url{https://par.nsf.gov/servlets/purl/10580612}.
\item Supplied literature reference: ``original OWR report,'' \url{https://doi.org/10.4171/OWR/2017/44}.
\item Literature assessment checked at \texttt{2026-08-21}: supplied dated curation, not independently verified. It labels the record partially solved, reports that the full Cohen--Macaulay log-concavity conjecture is disproved by the Niesi--Robbiano example, and describes standard graded Gorenstein-domain unimodality as open in general. The supplied background expressly requests independent verification of time-sensitive status and cited claims.
\item The registry supplies the label \texttt{corrected\_verified} and a statement-verification note reporting manual semantic, typographic, and LaTeX review of the clean statement. These are supplied editorial qualifications and do not establish mathematical correctness or independent source verification.
\end{itemize}
\end{document}
